% Ad Familiares 9.24 --- headnote
% ad-familiares-09-24, c. February 43 BC, Rome

\textit{Cicero to L. Papirius Paetus, written at Rome before the
middle of February 43 BC (Perseus: \textit{ante mcd. m. Febr.}). The
date places the letter in the thick of the Philippic period: the
First Philippic had been delivered in September 44, the Fourteenth
followed in late April 43, and the senatorial campaign that
eventually produced the disaster at Mutina is in full motion. This
is among the last surviving letters in Book 9 to Paetus, the
Campanian \textit{eques} and literary friend with whom Cicero had
exchanged the long, jocular dinner-letters of 46 and 45 BC.}

\textit{The letter has two registers, joined by a hinge. The first
section transacts business: Rufus, a friend of Paetus, has shown
concern for Cicero's safety; plots against Cicero at Aquinum and
Fabrateria have been thwarted in part because of Paetus's earlier
warnings, and so the recommendation Paetus has now sent twice is
already redundant from love. The middle section drops into the old
tone of teasing about dinners: Paetus has stopped going out, and
Cicero --- with Spurinna the haruspex as the deadpan medical
authority --- pretends to a grave Republican peril if Paetus does
not resume his old habits when the west wind (Favonius) blows. Then
the philosophical hinge, with the Greek pair \textit{symposia /
syndeipna} contrasted with the Latin \textit{convivia}: the Romans
were wiser to name the meal not by drink or food but by the sharing
of life. The closing turn pulls the register up sharply --- ``in
this care and administration if my life must be laid down, I shall
reckon that things have turned out splendidly for me'' --- and
gives the letter its valedictory weight.}
